\documentclass[runningheads]{llncs}

% ---------------------------------------------------------------
% Include basic ECCV package
\usepackage[review,year=2026,ID=2837]{eccv}

% ---------------------------------------------------------------
% Other packages
\usepackage{eccvabbrv}
\usepackage{graphicx}
\usepackage{booktabs}
\usepackage[numbers,sort&compress]{natbib}
\usepackage{amsmath}
\usepackage{amssymb}
\usepackage{multirow}
\usepackage[accsupp]{axessibility}
\usepackage[pagebackref,breaklinks,colorlinks,citecolor=eccvblue]{hyperref}
\usepackage{orcidlink}

\newcommand{\NA}{N/A}

\begin{document}

\title{Self-Evolving Unified Understanding and Generation Models through Unsupervised Self-Consistency: Supplementary Material}
\titlerunning{Self-Evolving UUG -- Supplementary}

\author{Anonymous ECCV Submission}
\authorrunning{Anonymous ECCV Submission}
\institute{Anonymous Institution}

\maketitle

\section{Protocol Comparison With Prior Work}
\label{sec:sup_protocol}

Table~\ref{tab:supervised_reference} summarizes supervision and optimization differences between our framework and commonly cited post-training baselines. This table is intended to qualify comparability, since many methods differ in reward supervision and architecture assumptions.

\begin{table*}[tb]
  \caption{Protocol-level comparison with prior post-training methods. "Extra supervision" indicates dependence on curated labels, external reward data, or supervised intermediate stages.}
  \label{tab:supervised_reference}
  \centering
  \small
  \resizebox{\linewidth}{!}{%
  \begin{tabular}{@{}lllll@{}}
    \toprule
    Method & Backbone setting in paper & Core optimization & Extra supervision & Notes on comparability \\
    \midrule
    UniRL~\cite{mao2025unirl} & Janus/Show-o family & Iterative SFT + GRPO & Yes (SFT/reward construction) & Single-family evaluation \\
    UniCorn~\cite{han2026unicorn} & Unified single-model setting & Role-based self-play reconstruction & Partial (judge/data reconstruction stage) & Limited cross-backbone evidence \\
    CoRL~\cite{jiang2025coreinforcement} & Janus-Pro-1B & Multi-stage RL & Yes (task-specific reward design) & Additional reward engineering \\
    UAE~\cite{yan2025understandinggeneration} & Qwen2.5-VL-3B & Unified-GRPO + staged reconstruction & Yes (supervised initialization stage) & Encoder--decoder coupling assumptions \\
    Internal Gap~\cite{han2025internalgap} & Janus-Pro-7B & SFT + DPO pipeline & Yes (curated preference-style data) & Understanding-guided data curation \\
    \midrule
    \textbf{Ours} & BLIP3o / BAGEL / VARGPT$_{1.1}$ & Alternating U/G self-evolving RL & \textbf{No} & Same framework across three generation paradigms \\
    \bottomrule
  \end{tabular}%
  }
\end{table*}


\section{Expanded Quantitative Results}
\label{sec:sup_expanded}

Table~\ref{tab:expanded_understanding} reports base vs. self-evolved results across all three backbones used in this work. We report only benchmarks that are available under each backbone's current evaluation pipeline; unavailable metrics are marked as \NA.

\begin{table*}[tb]
  \caption{Per-backbone base vs. self-evolved results across understanding and generation benchmarks.}
  \label{tab:expanded_understanding}
  \centering
  \small
  \resizebox{\linewidth}{!}{%
  \begin{tabular}{@{}llccccccccc@{}}
    \toprule
    Model & Setting & MMMU & MMBench & TextVQA & SEED & RWQA & MMVet & MME-P & MME-C & GenEval \\
    \midrule
    \multirow{3}{*}{BLIP3o-8B} & Base & 50.6 & 83.5 & 83.1 & 77.5 & 69.0 & 66.6 & 1682.6 & 647.1 & 0.84 \\
    & Ours & \textbf{52.8} & \textbf{86.1} & \textbf{85.2} & \textbf{79.4} & \textbf{70.9} & \textbf{68.7} & \textbf{1698.4} & \textbf{660.3} & \textbf{0.87} \\
    & $\Delta$ & +2.2 & +2.6 & +2.1 & +1.9 & +1.9 & +2.1 & +15.8 & +13.2 & +0.03 \\
    \midrule
    \multirow{3}{*}{BAGEL} & Base & 55.3 & 85.0 & \NA & \NA & \NA & 67.2 & 1687.0 & 701.0 & 0.82 \\
    & Ours & \textbf{58.8} & \textbf{87.1} & \NA & \NA & \NA & \textbf{69.5} & \textbf{1701.7} & \textbf{715.9} & \textbf{0.85} \\
    & $\Delta$ & +3.5 & +2.1 & \NA & \NA & \NA & +2.3 & +14.7 & +14.9 & +0.03 \\
    \midrule
    \multirow{3}{*}{VARGPT$_{1.1}$} & Base & 36.4 & 67.6 & 54.1 & 67.9 & 65.2 & 33.5 & 1488.8 & 348.6 & 0.33 \\
    & Ours & \textbf{39.4} & \textbf{70.3} & \textbf{56.9} & \textbf{71.0} & \textbf{68.8} & \textbf{36.2} & \textbf{1506.2} & \textbf{362.1} & \textbf{0.36} \\
    & $\Delta$ & +3.0 & +2.7 & +2.8 & +3.1 & +3.6 & +2.7 & +17.4 & +13.5 & +0.03 \\
    \bottomrule
  \end{tabular}%
  }
\end{table*}


\section{Implementation Details}
\label{sec:sup_implementation}

\noindent\textbf{Core optimization settings.}
All main runs use 10{,}000 steps, bfloat16 training, learning rate $1\times10^{-6}$, weight decay 0.01, gradient clip 1.0, and gradient accumulation 1. LoRA is enabled with rank 16, $\alpha=32$, dropout 0.05, and target modules: q\_proj, k\_proj, v\_proj, o\_proj, gate\_proj, up\_proj, down\_proj, and mm\_projector.

\noindent\textbf{Sampling and generation settings.}
Solver sampling uses 7 prompt perturbations. The proposer uses 3 candidates per step and 3 spot-check samples. Generation uses 3 candidates per prompt, image size 896$\times$896, 50 denoising/inference steps, and guidance scale 2.0 for BLIP3o-based runs.

\noindent\textbf{Reward and regularization.}
Generation reward weights are: QA fidelity 0.65, cycle-consistency 0.20, diversity 0.10, contradiction 0.20. Quality gates use minimum spec quality 0.35 and minimum 2 QA pairs per sample. KL control uses coefficient 0.01, target 0.02, adaptation rate 0.10, and bounds $[10^{-3},10^2]$.

\noindent\textbf{Cycle schedules.}
The main joint run uses a 3:2 understanding:generation cycle (E1). Understanding-only and generation-only ablations use 5:0 (E2) and 0:5 (E3), respectively. The imageless loop (E5) also uses 3:2, with replay-buffer-based generated-image training enabled (buffer size 500; generated mix ratio fixed to 1.0).

\noindent\textbf{Determinism and logging.}
All final scripts run with deterministic mode enabled, save checkpoints every 50 steps, and use seed 42 in the default configuration.


\bibliographystyle{splncs04}
\bibliography{main}

\end{document}
